% Cap.~1 Fig.~1.2 --- throughput Monte Carlo as a black box
\begin{figure}[htbp]
\centering
\resizebox{\textwidth}{!}{%
\begin{tikzpicture}[>=Stealth]

% ===== Left: black box =====
\node[
  draw=red!35, rounded corners=4pt, fill=red!4,
  minimum width=6.8cm, minimum height=5.6cm
] (LP) at (0,0) {};

\node[font=\bfseries\small, text=red!65!black]
  at (0,2.3) {Black-box approach};

\node[font=\small] (cnt) at (-2.35,0.6) {count};
\node[
  fill=black, text=white, font=\bfseries\small, align=center,
  minimum width=3.0cm, minimum height=1.3cm, rounded corners=1pt
] (bb) at (0,0.6) {Mean of\\completions};
\node[font=\small] (fc) at (2.4,0.6) {forecast};
\draw[->, ultra thick, red!60!black] (cnt) -- (bb);
\draw[->, ultra thick, red!60!black] (bb) -- (fc);

\node[font=\footnotesize, align=center, text=black!70]
  at (0,-1.4)
  {Sees only the \textbf{how many}.\\[3pt]Ignores the \textbf{how}.};

% ===== Right: three separated cards =====
% Card 1
\node[circle, fill=red!70!black, text=white, font=\bfseries\scriptsize,
  minimum size=0.45cm, inner sep=0pt] (n1) at (4.15,2.05) {1};
\node[
  draw=black!25, rounded corners=3pt, fill=white, anchor=west,
  text width=7.0cm, align=left, inner sep=6pt
] at (4.55,2.05)
  {\textbf{Assumes stability}\\[2pt]
   \footnotesize Treats every week as exchangeable. Real processes change:
   rework, bottlenecks, seasonality.};

% Card 2
\node[circle, fill=red!70!black, text=white, font=\bfseries\scriptsize,
  minimum size=0.45cm, inner sep=0pt] (n2) at (4.15,0.15) {2};
\node[
  draw=black!25, rounded corners=3pt, fill=white, anchor=west,
  text width=7.0cm, align=left, inner sep=6pt
] at (4.55,0.15)
  {\textbf{Loses structure}\\[2pt]
   \footnotesize A case in Review is treated like one in To Do; state
   position is discarded.};

% Card 3
\node[circle, fill=red!70!black, text=white, font=\bfseries\scriptsize,
  minimum size=0.45cm, inner sep=0pt] (n3) at (4.15,-1.75) {3};
\node[
  draw=black!25, rounded corners=3pt, fill=white, anchor=west,
  text width=7.0cm, align=left, inner sep=6pt
] at (4.55,-1.75)
  {\textbf{Does not explain}\\[2pt]
   \footnotesize Yields a number, not a cause---no actionable bottleneck
   diagnosis.};

\end{tikzpicture}%
}
\caption{Structural limits of throughput-based Monte Carlo (market practice):
weekly completion counts are sampled to project backlog burn-down, but the
method never opens the process.
Throughput is how many tasks finish per week; the black box sees the
\emph{count}, not the \emph{path}.
At root it is a sophisticated average that assumes a generative process
rarely present in practice.}
\label{fig:cap1-throughput-blackbox}
\end{figure}
